% Source pp. 227-238 (Lectures19-25.pdf).
\section{Nonlinear Sigma Model at \texorpdfstring{$d=2+\epsilon$}{d=2+epsilon}}

Now I would like to discuss how the RG transformation acts on the action of the
$\sigma$-model at $d=2$. In this case the coupling constant $g$ in
\[
\mathcal A=\int d^2x\,\frac1{2g}\,(\partial\mathbf n)^2
\]
is dimensionless. I assume that $g$ is small, so that the fluctuations of
$\mathbf n(x)$ are small, since at small $g$ the exponential factor
$\exp(-\mathcal A)$ strongly suppresses configurations having large gradients
$\partial\mathbf n$. I remind you that the field $\mathbf n(x)$ satisfies the
non-linear constraint $\mathbf n^2(x)=1$.

In principle, it is not difficult to develop renormalized perturbation theory in
the coupling constant $g$, derive the CS equation and find the series expansion of
associated flow function $\beta(g)$. This calculation is displayed in \S\,13.3 of
PS. I will discuss the derivation based explicitly on the idea of integrating out
the fast modes (also described in \S\,13.3 of PS), because it has some novel
features as compared to what we have done before.

First, we introduce cutoff,
\[
\mathbf n(x)=\int_{|k|<\Lambda}\mathbf n_k\,e^{ikx}\,\frac{d^2k}{(2\pi)^2}\,.
\]
Note that now the modes $\mathbf n_k$ are not all independent, but must obey
complicated constraint $\int_q\mathbf n_q\mathbf n_{k-q}=\delta^{(2)}(k)$, because
of the $\sigma$-model constraint $\mathbf n^2(x)=1$.

Next, we want to split $\mathbf n(x)$ into its ``slow'' and ``fast'' components.
The above nonlinear constraint makes this task less straightforward than we had it
before in usual scalar theories. Physically transparent way to make the splitting
is as follows. We define the ``slow'' component $\mathbf n_0(x)$ as the field of
unit length,
\begin{equation}\label{eq:n0unit}
\mathbf n_0^2(x)=1\,,
\end{equation}
whose Fourier transform is restricted to $|k|<\Lambda/L$. To make further
arguments transparent, I will assume that $L\gg1$; we will see that in this case
the dominating contribution to the RG transformed action comes from the
integration over the fast modes with $\Lambda/L\ll|k|\ll\Lambda$. More precisely,
I will assume that $L\gg1$, but $g$ is so small that
\[
g\log L\ll1\,.
\]
At each spatial point $x$, one can choose a system of vectors $\mathbf e_i(x)$,
$i=1,2,\dots,N-1$, such that $N$ vectors $(\mathbf n_0(x),\mathbf e_i(x))$ form a
basis. It is convenient to take orthonormal basis,
\begin{equation}\label{eq:ebasis}
\bigl(\mathbf n_0(x)\,\mathbf e_i(x)\bigr)=0\,,\qquad\bigl(\mathbf e_i(x)\,\mathbf e_k(x)\bigr)=\delta_{ik}\,.
\end{equation}
Since $\mathbf e_i(x)$ only need to be orthogonal to the ``slow'' field
$\mathbf n_0(x)$, one can always choose them to be ``slow'' fields as well,
containing only Fourier modes with $|k|$ of the order of $\Lambda/L$ or smaller.
The full field $\mathbf n(x)$ can be written as
\begin{equation}\label{eq:nsplit}
\mathbf n(x)=\mathbf n_0(x)\,\sqrt{1-(\phi_i)^2(x)}+\mathbf e^i(x)\,\phi_i(x)\,.
\end{equation}
This expression automatically satisfies the constraint $\mathbf n^2(x)=1$. $N-1$
fields $\phi_i(x)$ represent the fast modes; they contain Fourier components with
$|k|$ between $\Lambda/L$ and $\Lambda$. It is important that when $g$ is small we
can assume $\phi_i(x)$ are small.

In order to compute the action, we need spatial derivatives of $\mathbf n(x)$.
These involve the derivatives of the vectors $\mathbf n_0(x)$ and $\mathbf e_i(x)$,
as well as of the fields $\phi_i$. One can expand them in the basic vectors as
\begin{align}
\partial_\mu\mathbf n_0(x)&=B_\mu^i(x)\,\mathbf e_i(x)\,,\label{eq:dn0}\\
\partial_\mu\mathbf e_i(x)&=A_\mu^{ik}(x)\,\mathbf e_k(x)-B_\mu^i(x)\,\mathbf n_0(x)\,,\label{eq:de}
\end{align}
where
\begin{align}
B_\mu^i(x)&=(\mathbf e_i\,\partial_\mu\mathbf n_0)=-(\mathbf n_0\,\partial_\mu\mathbf e_i)\,,\label{eq:Bdef}\\
A_\mu^{ik}(x)&=(\mathbf e_k\,\partial_\mu\mathbf e_i)=-(\mathbf e_i\,\partial_\mu\mathbf e_k)\,.\label{eq:Adef}
\end{align}
Note that the fields $A_\mu^{ik}$ and $B_\mu^i$ are ``slow''. In terms of these
variables the action is
\[
\mathcal A=\frac1{2g}\int d^2x\,\Bigl[\Bigl(\partial_\mu\sqrt{1-\phi^2}-B_\mu^i\partial_\mu\phi_i\Bigr)^2
+\Bigl(\partial_\mu\phi_i-A_\mu^{ik}\phi_k+B_\mu^i\sqrt{1-\phi^2}\Bigr)^2\Bigr]\,.
\]
For small $g$ one can neglect all terms beyond the terms quadratic in $\phi$,
\[
\mathcal A=\mathcal A[\mathbf n_0]+\tilde{\mathcal A}[\mathbf n_0,\phi_i]\,,
\]
where the first term comes from
\[
\frac1{2g}\int d^2x\,B_\mu^iB_\mu^i=\frac1{2g}\int d^2x\,(\partial_\mu\mathbf n_0)^2=\mathcal A[\mathbf n_0]\,,
\]
and
\[
\tilde{\mathcal A}=\frac1{2g}\int d^2x\,\Bigl[\bigl(\partial_\mu\phi_i-A_\mu^{ik}\phi_k\bigr)^2+B_\mu^iB_\mu^k\bigl(\phi_i\phi_k-\delta_{ik}\phi^2\bigr)\Bigr]\,.
\]
Our goal now is to integrate out the fast modes $\phi$. We treat the term
$(1/2g)\int(\partial\phi)^2$ as the unperturbed action, yielding the propagator
for the $\phi$-field
\[
i\ \dia{\draw[prop] (-1.6,0) -- (1.6,0);}\ k=\frac{g\,\delta_{ik}}{p^2}\,.
\]
The remaining terms are treated as perturbations, giving rise to vertices of
three types:
\begin{align*}
\diac[0.45]{\draw[prop,densely dashed] (-2,0) -- (0,0); \draw[prop] (0,0) -- (1.3,1.3) node[above right] {\small$j$}; \draw[prop] (0,0) -- (1.3,-1.3) node[below right] {\small$k$};
  \node[vtx,minimum size=4.5pt] at (0,0) {}; \node[left] at (-2,0) {\small$A$};
  \draw[arr] (0.95,1.35) -- (0.45,0.85); \node[font=\tiny] at (0.4,1.3) {$p_1$};
  \draw[arr] (0.95,-1.35) -- (0.45,-0.85); \node[font=\tiny] at (0.4,-1.3) {$p_2$};}
&=\frac{i\,A_\mu^{jk}}g\,(p_1-p_2)_\mu\,,\\[4pt]
\diac[0.45]{\draw[prop,densely dashed] (-2,0) -- (0,0); \draw[prop] (0,0) -- (1.3,1.3) node[above right] {\small$j$}; \draw[prop] (0,0) -- (1.3,-1.3) node[below right] {\small$k$};
  \node[vtx,minimum size=4.5pt] at (0,0) {}; \node[left] at (-2,0) {\small$A^2$};}
&=-\frac{A_\mu^{ij}A_\mu^{ik}}g\,,\\[4pt]
\diac[0.45]{\draw[prop,densely dashed] (-2,0) -- (0,0); \draw[prop] (0,0) -- (1.3,1.3) node[above right] {\small$j$}; \draw[prop] (0,0) -- (1.3,-1.3) node[below right] {\small$k$};
  \node[vtx,minimum size=4.5pt] at (0,0) {}; \node[left] at (-2,0) {\small$B^2$};}
&=\frac{\delta_{jk}\,B_\mu^iB_\mu^i-B_\mu^jB_\mu^k}g\,.
\end{align*}
As usual
\[
\mathcal A_1=\mathcal A-\text{sum of all connected diagrams}\,,
\]
with the dashed lines taken as the external legs.

The leading order in $g$ comes from the one-loop diagrams. The simplest (and the
most important) of those is
\begin{align*}
\dia[0.45]{\draw[prop,densely dashed] (-2,0) -- (0,0); \draw[prop] (0.9,0) circle (0.9); \node[vtx,minimum size=4.5pt] at (0,0) {}; \node[left] at (-2,0) {\small$B^2$};}
&=\frac{\delta_{jk}\,B_\mu^iB_\mu^i-B_\mu^jB_\mu^k}{2g}\int\frac{d^2p}{(2\pi)^2}\,\frac{g\,\delta_{jk}}{p^2}\\
&=\frac{(N-2)\,(B_\mu^i)^2}2\,\frac1{4\pi}\int_{\Lambda^2/L^2}^{\Lambda^2}\frac{dp^2}{p^2}=\frac{N-2}{4\pi}\,(\partial_\mu\mathbf n_0)^2\,\log L\,.
\end{align*}
Thus, this diagram brings in a term identical in the form to the original action,
i.e.\ it renormalizes the coupling parameter $g$.

The diagram
\[
\dia[0.45]{\draw[prop,densely dashed] (-2.6,0) -- (-1.2,0) (1.2,0) -- (2.6,0); \draw[prop] (0,0) ellipse (1.2 and 0.75);
  \node[vtx] at (-1.2,0) {}; \node[vtx] at (1.2,0) {}; \node[left] at (-2.6,0) {\small$B^2$}; \node[right] at (2.6,0) {\small$B^2$};}
\]
generates term $\sim B^4$ (with the coefficient $\sim1/\Lambda^2$); the field
$B^4$ has dimension $4>d$ and hence it is an irrelevant field; it can be neglected
in the leading approximation. Likewise, all the diagrams of the form
\[
\diac[0.45]{\draw[prop] (0,0) circle (1.1);
  \foreach \a in {180,120,60,0} {\node[vtx] at (\a:1.1) {}; \draw[prop,densely dashed] (\a:1.1) -- (\a:2.2); \node at (\a:2.6) {\scriptsize$B^2$};}
  \draw[densely dotted] (200:1.5) arc[start angle=200,end angle=340,radius=1.5];}
\]
can be ignored, for the same reason.

Let us take a look at the diagrams with $A$-vertices. Consider the following ones
\[
\dia[0.45]{\draw[prop,densely dashed] (-2,0) -- (0,0); \draw[prop] (0.9,0) circle (0.9); \node[vtx,minimum size=4.5pt] at (0,0) {}; \node[left] at (-2,0) {\small$A^2$};}\,,\qquad
\dia[0.45]{\draw[prop,densely dashed] (-2.6,0) -- (-1.2,0) (1.2,0) -- (2.6,0); \draw[prop] (0,0) ellipse (1.2 and 0.75);
  \node[vtx] at (-1.2,0) {}; \node[vtx] at (1.2,0) {}; \node[left] at (-2.6,0) {\small$A$}; \node[right] at (2.6,0) {\small$A$};}\,.
\]
Both these diagrams show logarithmic behavior,
\[
\int\frac{d^2k}{k^2}\sim\log L\,,\qquad\text{and}\qquad\int\frac{d^2k\,k^\mu k^\nu}{k^2\,k^2}\sim\log L\,.
\]
However, it is not difficult to show that these two logarithmic contributions in
fact cancel each other. One can do it by a direct calculation, but it is easy to
prove this statement on the basis of certain symmetry.

The $\log L$ of each of these two diagrams is of the form
\[
\log(L)\,A_\mu^{ij}A_\mu^{ij}\,,
\]
where, as before,
\[
A_\mu^{ij}=\mathbf e^j\,\partial_\mu\mathbf e^i\,.
\]
Recall that $\mathbf e^i(x)$ are arbitrary vectors tangent to $S^{N-1}$ at
$\mathbf n_0(x)$. They satisfy
\[
\mathbf e^i\mathbf n_0=0\,,\qquad\mathbf e^i\mathbf e^j=\delta^{ij}\,,
\]
but otherwise arbitrary. One can always choose them differently,
\[
\tilde{\mathbf e}^i(x)=\Omega_{ij}(x)\,\mathbf e^j(x)\,,
\]
where $\Omega_{ij}(x)$ is orthogonal $(N-1)\times(N-1)$ matrix,
\[
\Omega(x)\,\Omega^t(x)=I\,.
\]
$\Omega(x)$ is ``slow'' (for $\mathbf e^i(x)$ were assumed to be slow), but
otherwise arbitrary function of $x$. Under these local orthogonal transformations,
$A_\mu^{ij}(x)$ transforms as an $O(N-1)$ gauge potential,
\[
\tilde A_\mu(x)=\Omega(x)\,A_\mu(x)\,\Omega^t(x)+\Omega(x)\,\partial_\mu\Omega^t(x)\,,
\]
where $A_\mu(x)$ stands for $(N-1)\times(N-1)$ anti-symmetric matrix
$A_\mu^{ij}(x)$.

As this gauge transformation describes just the freedom in the choice of
variables, the action $\mathcal A_1[\mathbf n_0]$ must be gauge invariant. It is
easy to check that the combination
\[
A_\mu^{ij}A_\mu^{ij}=\tr\bigl(A_\mu A_\mu^t\bigr)
\]
is not gauge invariant, and thus it can not appear in $\mathcal A_1$. The simplest
gauge invariant local combination of course is
\[
\tr\bigl(F_{\mu\nu}F_{\mu\nu}\bigr)\,,\qquad\text{with}\qquad
F_{\mu\nu}=\partial_\mu A_\nu-\partial_\nu A_\mu+[A_\mu,A_\nu]\,.
\]
However it involves extra derivative, and thus corresponds to an irrelevant term.

Therefore, in this approximation the only interesting contribution comes from the
above loop diagram with single $B^2$ vertex. We have
\begin{equation}\label{eq:A1sigma}
\mathcal A_1=\int\Bigl(\frac1{2g}-\frac{N-2}{4\pi}\,\log L\Bigr)\,(\partial_\mu\mathbf n_0)^2\,d^2x\,.
\end{equation}
Thus, the step 1 of the RG transformation leads to
\[
\frac1g\ \to\ \frac1{g_1}=\frac1g-\frac{N-2}{2\pi}\,\log L\,.
\]
The step 2 is nearly trivial. If one changes
\[
\mathbf n_0(x)=\mathbf n(x/L)
\]
the action does not change. Therefore
\[
RG_l:\qquad\frac1g\ \to\ \frac1g-\frac{N-2}{2\pi}\,l\,.
\]
In other words, the beta-function is
\[
\beta(g)=-\frac{N-2}{2\pi}\,g^2+O(g^3)\,,
\]
and for small $g$ the RG flow equation
\[
\frac{dg}{dl}=\frac{N-2}{2\pi}\,g^2
\]
shows that $g(l)$ \emph{increases} towards large distances; at $d=2$ the Goldstone
fixed point becomes unstable.

The above flow equation determines (in the leading order in $g$) the behavior of
the action under the RG transformation. One might be interested in the behavior
of correlation functions of, say, the field $\mathbf n(x)$. To that end one needs
then to find the RG function $\gamma_{\mathbf n}(g)$ in the CS equation
\[
\Bigl(x\frac\partial{\partial x}+n\,\gamma_{\mathbf n}(g)+\beta(g)\,\frac\partial{\partial g}\Bigr)\bigl\langle\mathbf n(x_1)\cdots\mathbf n(x_n)\bigr\rangle_c=0\,.
\]
It is defined in terms of the action of the RG transformation on the field
$\mathbf n(x)$.

Consider any correlation function involving $\mathbf n(x)$,
\[
\bigl\langle\mathbf n(x)\cdots\bigr\rangle\,,
\]
and assume again that the momentum $p$ associated with $x$ is $\ll\Lambda/L$.
Then, splitting as before
\[
\mathbf n=\mathbf n_0\,\sqrt{1-\phi^2}+\mathbf e^i\,\phi_i\simeq\mathbf n_0\,\Bigl(1-\frac12\,\phi^2\Bigr)+\mathbf e^i\,\phi_i
\]
one needs to find expectation value of this combination over the ensemble of the
fast modes $\phi$. The last term can be ignored (its expectation value is zero),
and we have
\[
\mathbf n\ \to\ \mathbf n_0\,\Bigl(1-\frac12\,\langle\phi_i^2\rangle\Bigr):=Z_{\mathbf n}^{-\frac12}(L)\,\mathbf n_0\,,
\qquad\text{where}\qquad
Z_{\mathbf n}^{-\frac12}(L)=1-\frac12\,\langle\phi_i^2\rangle=1-\frac{N-1}2\,\frac g{2\pi}\,l\,.
\]
Since at $d=2$ $\mathbf n$ has zero canonical dimension, we then have
\[
\gamma_{\mathbf n}(g)=D_{\mathbf n}(g)=-\frac d{dl}\,Z_{\mathbf n}^{-\frac12}(l)\Big|_{l=0}=\frac{N-1}{4\pi}\,g\,.
\]
In the above calculation we have assumed that $gl\ll1$; that is why the
higher-order terms in $\phi$ were neglected. But, as usual, we can extend it to
larger $l$. Generally, we have
\[
RG_l\,\mathbf n=Z_{\mathbf n}^{-\frac12}(l)\,\mathbf n\,,
\]
where $Z_{\mathbf n}^{-\frac12}(l)$ is a solution of the equation
\[
-\frac d{dl}\,Z_{\mathbf n}^{-\frac12}(l)=\gamma_{\mathbf n}\bigl(g(l)\bigr)\,Z_{\mathbf n}^{-\frac12}(l)\,,
\]
and $g(l)$ is the integral curve of the equation
\[
\frac{dg(l)}{dl}=-\beta\bigl(g(l)\bigr)\,.
\]
In our case
\[
\frac1{g(l)}=\frac1{g(0)}-\frac{N-2}{2\pi}\,l\,.
\]
Recall that the ``flow'' of $Z_{\mathbf n}^{-\frac12}(l)$ determines the scale
dependence of the correlation functions
\[
\bigl\langle\mathbf n(x_1/L)\cdots\mathbf n(x_n/L)\bigr\rangle_{g(l)}=Z_{\mathbf n}^{-\frac n2}(l)\,\bigl\langle\mathbf n(x_1)\cdots\mathbf n(x_n)\bigr\rangle_{g(0)}\,,
\]
where $l=\log L$.

\begin{exercise}
For the $O(N)$ nonlinear sigma model,

I. At $d=2$
\begin{enumerate}[label=\alph*)]
\item Assuming $g(l)\ll1$, find $Z_{\mathbf n}(l)$.
\item Derive the following $|x|\to0$ asymptotic of the two-point correlation
function
\[
\bigl\langle\mathbf n(x)\mathbf n(0)\bigr\rangle\sim\Bigl[\log\Bigl(\frac1{x^2}\Bigr)\Bigr]^{\frac{N-1}{N-2}}\,.
\]
Explain how this result is in agreement with the exact solution of the model at
$N\to\infty$ (PS, \S13.3, pages 463--465).
\end{enumerate}
II. Find the flow equation for $g(l)$ and $Z_{\mathbf n}(l)$ at $d=2+\epsilon$,
with $\epsilon\ll1$. Again, determine the $|x|\to0$ behavior of the two-point
correlation function.
\end{exercise}

Coming back to the flow equation ($d=2$)
\[
\frac{dg}{dl}=\frac{N-2}{2\pi}\,g^2+\dots
\]
one observes that (at $N>2$) the flow of the coupling $g(l)$ looks like this:
\[
\diac[0.6]{\draw[line width=0.8pt] (0,0) -- (6.5,0) node[right] {$g$}; \draw[arr] (2.5,0) -- (3.6,0);
  \fill (0,0) circle (2.4pt) node[above=2pt] {$0$};}
\]
At large distances $g$ increases, and the interaction becomes strong. Our weak
coupling analysis is no longer valid (in particular the higher order terms in
$\beta(g)$ can not be neglected). On the other hand, if we follow the flow of
$g(l)$ backward, to $l\to-\infty$,
\[
g(l)=\Bigl(\Const-\frac{N-2}{2\pi}\,l\Bigr)^{-1}\ \to\ \frac{2\pi}{N-2}\,(-l)^{-1}\,,
\]
the interaction of the Goldstone modes becomes weaker, and the short-distance
behavior is controlled by a free theory. For this reason, this type of RG flow is
called \emph{asymptotic freedom}. Note that the situation is opposite to that in
$\varphi^4$ theory at $d=4$, where instead $\lambda\to0$ at long scales, whereas
perturbation theory breaks down at short distances.

The asymptotic freedom allows one to establish existence of local field theory;
indeed it shows that at short scales the fluctuations are suppressed, and the
cutoff can be removed with no violation of locality. On the other hand, at longer
scales the fluctuations become stronger. It is not possible to describe these
large scale fluctuations by perturbation theory. Nonetheless, one can make
intuitive picture based on this pattern of the RG flow. If one concentrates
attention on very short scales, the vectors $\mathbf n(x)$ are nearly parallel,
because at $l\to-\infty$ $g(l)$ is small, thus suppressing the gradients of
$\mathbf n(x)$. However, if one looks at larger scales, the directions of
$\mathbf n(x)$ at the points separated by distance $\sim R_c$ are much less
correlated, and if the separation is even much larger, there is no correlation
(figure~\ref{fig:sigmascales}).
\begin{figure}[ht]
\centering\figfile{sigmascales}
\caption{At short scales the vectors $\mathbf n(x)$ are nearly parallel (left).
At scales much larger than $R_c$ their directions are uncorrelated (right).}
\label{fig:sigmascales}
\end{figure}

This picture suggests that at $d=2$ the phase with spontaneously broken $O(N)$
symmetry does not exist. In other words, $d=2$ magnet has no ferromagnetic phase
with spontaneous magnetization. Let me stress that the perturbative analysis of
the RG flow does not prove this statement, which requires much more involved
arguments. Note also that this picture is expected to apply only if
$N>2$.\footnote{At $N=2$ the $g^2$ term in the $\beta$-function vanishes. In fact,
it is easy to see that all higher terms vanish as well. At $N=2$ the direction
$\mathbf n(x)$ can be parameterized by a single angle $\theta(x)$,
$\mathbf n(x)=\bigl(\cos\theta(x),\sin\theta(x)\bigr)$, and formally
$\mathcal A=\frac1{2g}\int(\partial_\mu\theta)^2\,d^2x$, which is just the action
of free massless theory. The action of RG on a free theory is trivial. However,
beyond the perturbation theory, the $N=2$ sigma model has more complex structure.
This is because the field $\theta(x)$ is an angular variable,
$\theta(x)\sim\theta(x)+2\pi N(x)$, where $N(x)$ is integer-valued. As the result,
there are topologically nontrivial field configurations (``vortices'', or
``instantons''). Their interaction leads to remarkable phase transition at finite
$g=g_c$, the so called Berezinsky-Kosterlitz-Thouless transition (See e.g.\
Kosterlitz, Thouless, ``Ordering, metastability and phase transitions in
two-dimensional systems'', Journal of Physics C: Solid State Physics, Vol.~6 pages
1181--1203 (1973)).}

What is different at $d>2$? This question can be studied in details in formal case
of $d=2+\epsilon$ with $\epsilon\ll1$. In this case we can repeat our calculation
of the RG flow at small $g$. The step 1 (elimination of the fast part) yields
essentially the same result,
\begin{equation}\label{eq:A1sigmaeps}
\mathcal A_1=\int\Bigl(\frac1{2g}-\frac{N-2}{4\pi}\,\Lambda^\epsilon\,\log L\Bigr)\,(\partial_\mu\mathbf n_0)^2\,d^dx\,.
\end{equation}
However, the step 2 (scale transformation) is no longer exactly trivial. The change
of variables
\[
\mathbf n_0(x)=\mathbf n(x/L)
\]
leads to
\[
\int(\partial_\mu\mathbf n_0)^2\,d^dx=L^{d-2}\int(\partial_\mu\mathbf n)^2\,d^dx\,.
\]
Therefore
\[
\mathcal A'=\int L^\epsilon\,\Bigl(\frac1{2g}-\frac{N-2}{4\pi}\,\Lambda^\epsilon\,\log L\Bigr)\,(\partial_\mu\mathbf n_0)^2\,d^dx\,.
\]
Since
\[
L^\epsilon\approx1+\epsilon\log L
\]
one finds that under the RG transformation
\[
\frac1g\ \to\ \frac1g+\Bigl(\frac\epsilon g-\frac{N-2}{2\pi}\Bigr)\,\log L\,,
\]
where I have ignored the factor $\Lambda^\epsilon\approx1$.\footnote{One can get
rid of this factor by changing to a dimensionless coupling $g\to\Lambda^{-\epsilon}g$.}
The associated RG flow equation is
\[
\frac{dg}{dl}=-\epsilon\,g+\frac{N-2}{2\pi}\,g^2=\frac{N-2}{2\pi}\,g\,(g-g_*)\,,
\qquad\text{where}\qquad
g_*=\frac{2\pi\,\epsilon}{N-2}\,.
\]
Due to the small linear term, the flow now looks different at small $g\sim\epsilon$,
\[
\diac[0.6]{\draw[line width=0.8pt] (0,0) -- (6.5,0) node[right] {$g$};
  \draw[arr] (2.0,0) -- (0.9,0); \draw[arr] (3.8,0) -- (5.0,0);
  \node at (0,0) {$*$}; \node[above=2pt] at (0,0) {GS}; \node at (2.9,0) {$*$}; \node[above=2pt] at (2.9,0) {$g_*$};}
\]
Nontrivial new fixed point at $g=g_*$ appears, and if $g<g_*$ the theory flows
towards the Goldstone fixed point at large scale.

It is natural to identify the new fixed point $g_*$ with the Wilson-Fisher fixed
point at $d=2+\epsilon$. Indeed, general pattern of RG flow in $O(N)$ $\varphi^4$
theory suggested by our previous analysis was as shown in
figure~\ref{fig:wfgs}.
\begin{figure}[ht]
\centering\figfile{wfgs}
\caption{The flow from $\mathcal A_{*,WF}$ to the Goldstone point
$\mathcal A_{*,GS}$; the deviation $g-g_*$ is the sigma-model coordinate along it.}
\label{fig:wfgs}
\end{figure}

Here $M^2$ and $\lambda$ are two couplings in $\varphi^4$ theory. Closely below
$d=4$ the Wilson-Fisher fixed point appears close to $\mathcal A_{*,G}$, while the
Goldstone fixed point is located far away from it. When $d$ goes down from 4
$\mathcal A_{*,G}$ departs from $\mathcal A_{*,WF}$, but $\mathcal A_{*,GS}$ gets
closer (figure~\ref{fig:wfgs2}).
\begin{figure}[ht]
\centering\figfile{wfgs2}
\caption{The same flow at $d=3.99$ (left) and $d=2.01$ (right). As $d$ decreases
the Wilson-Fisher point moves away from the Gaussian point and towards the
Goldstone point.}
\label{fig:wfgs2}
\end{figure}

We can identify the coordinate along the trajectory going from
$\mathcal A_{*,WF}$ to $\mathcal A_{*,GS}$ with the coupling $g$ of the
non-linear sigma model, with the origin $g=0$ at $\mathcal A_{*,GS}$; then
$\mathcal A_{*,WF}$ appears at some $g_*$. At $d=2+\epsilon$ with small $\epsilon$
$\mathcal A_{*,GS}$ approaches $\mathcal A_{*,WF}$ ($g_*\sim\epsilon$). At $d=2$ the
fixed points $\mathcal A_{*,GS}$ and $\mathcal A_{*,WF}$ merge, and the flow of $g$
(at positive $g$) becomes marginally relevant. This pattern is confirmed by exact
solution of the theory at $N\to\infty$ (partly described in \S13.3 of PS).

\begin{hwproblem}{9}\label{hw:9}
For the $O(N)$ sigma model,
\[
\mathcal A[\mathbf n(x)]=\int d^2x\,\frac1{2g}\,\bigl(\partial\mathbf n(x)\bigr)^2\,.
\]
Assume $N>2$.

I. At $d=2$
\begin{hwparts}
\item Assuming $g(l)\ll1$, find $Z_{\mathbf n}(l)$ (see Section~22 for
definition).
\item Determine the $|x|\to0$ behaviour of the two-point correlation function
$\langle\mathbf n(x)\mathbf n(0)\rangle$.
\item \emph{Optional.} Explain how the result is in agreement with the exact
solution of the model at $N\to\infty$ (PS, \S13.3, pages 463--465).
\end{hwparts}
\emph{Direction:} Find $Z_{\mathbf n}(l)$ from
$-\frac d{dl}Z_{\mathbf n}^{-1/2}(l)=\gamma_{\mathbf n}(g(l))\,Z_{\mathbf n}^{-1/2}(l)$.
Use the equation
\[
\bigl\langle\mathbf n(x_1/L)\,\mathbf n(x_2/L)\bigr\rangle_{g(l)}=Z_{\mathbf n}^{-1}(l)\,\bigl\langle\mathbf n(x_1)\,\mathbf n(x_2)\bigr\rangle_{g(0)}\,,
\]
and remember that this correlation function depends on the coordinates through the
separation $|x_1-x_2|$ only.

II. At $d=2+\epsilon$ with $\epsilon\ll1$, find the flow equation for
$Z_{\mathbf n}(l)$. Again, determine the $|x|\to0$ behavior of the two-point
correlation function.
\end{hwproblem}
